\documentclass[11pt]{article}
\usepackage{url}
\usepackage{alltt}
\usepackage{bm}
\usepackage{bbm}
\linespread{1}
\textwidth 6.5in
\oddsidemargin 0.in
\addtolength{\topmargin}{-1in}
\addtolength{\textheight}{2in}

\usepackage{amsmath}
\usepackage{amssymb}
\usepackage{hyperref}

\begin{document}


\begin{center}
\Large
STA 711 Exam 2 Review\\
\normalsize
\vspace{5mm}
\end{center} 

\noindent \textbf{Information:} The second exam will focus on the material from homeworks 5 - 7, though some material from earlier assignments (such as maximum likelihood estimation and convergence results) may also appear. Here is a rough list of topics for the second exam, though I do not claim it is exhaustive:

\begin{itemize}
\item Asymptotic results, including: applications of the WLLN and CLT, the asymptotic normality of the MLE, asymptotic behavior of the MLE under mis-specification, delta method, asymptotic relative efficiency, and M-estimation and the sandwich variance
\item Estimation and comparing estimators: MLE, method of moments, comparing estimators with their MSE, unbiased estimators, and the Cramer-Rao lower bound and best unbiased estimators
\item Multivariate normal distributions
\end{itemize}

\noindent The practice questions below are not completely comprehensive, but will give you a sense of the kinds of questions I could ask on the exam.\\

\noindent \textbf{Resources and rules:} 

\begin{itemize}
\item You will be provided with a sheet of common distributions, including their probability function, mean, and variance. You may use any of the facts on this sheet without proof, unless specifically asked to derive them

\item You will also be provided with a partial list of results from class and homework assignments (see the next page); you may use these results without proof, unless asked to specifically to derive one of these results or otherwise stated.

\item Finally, any other results from class or homework may also be used without proof, unless asked to specifically to derive one of these results or otherwise stated.
\end{itemize}

\newpage

\section*{Provided results}

\subsection*{Quantile estimation}

Let $Y_1,...,Y_n \overset{iid}{\sim} F$, for some continuous $F$ with pdf $f$, and let $\theta_0$ be the $p$th quantile of $F$, and $\widehat{\theta}$ the sample $p$th quantile. Then,
$$\sqrt{n}(\widehat{\theta} - \theta_0) \overset{d}{\to} N\left(0, \frac{p(1-p)}{f^2(\theta_0)} \right)$$
(provided that $f(\theta_0) > 0$).

\subsection*{M-estimation (including extension)}

Let $Y_1,...,Y_n$ be iid samples from some distribution $g$. ($Y_i$ could be a scalar or vector random variable; and in the regression setting, we typically have $Y_i = (\mathbf{x}_i, Y_i)$). Let $\psi(Y, \boldsymbol{\theta})$ be a function which depends on a parameter $\boldsymbol{\theta} \in \mathbb{R}^d$.\\

\noindent Let $\widehat{\boldsymbol{\theta}}$ solve
$$\sum \limits_{i=1}^n \psi(Y_i, \boldsymbol{\theta}) = \boldsymbol{c}_n$$
where $\boldsymbol{c}_n/\sqrt{n} \overset{p}{\to} \boldsymbol{0}$ as $n \to \infty$ (in the simplest case, we just have $\boldsymbol{c}_n \equiv \boldsymbol{0}$). And let $\boldsymbol{\theta}_0$ solve
$$\mathbb{E}_g[\psi(Y_i, \boldsymbol{\theta})] = \boldsymbol{0}$$
Then, under regularity conditions (which you may assume hold unless otherwise specified), $\widehat{\boldsymbol{\theta}} \overset{p}{\to} \boldsymbol{\theta}_0$, and
$$\sqrt{n}(\widehat{\boldsymbol{\theta}} - \boldsymbol{\theta}_0) \overset{d}{\to} N(\boldsymbol{0}, \mathbf{S}(\boldsymbol{\theta}_0)) \hspace{1cm} \mathbf{S}(\boldsymbol{\theta}_0) = \mathbf{A}^{-1} \mathbf{B} \mathbf{A}^{-1}$$
where $\mathbf{B} = \mathbb{E}_g[\psi(Y_i, \boldsymbol{\theta}_0) \psi(Y_i, \boldsymbol{\theta}_0)^T]$. The definition of $\mathbf{A}$ depends on whether $\psi$ is differentiable. If $\psi$ is differentiable, then
$$\mathbf{A} = -\mathbb{E}_g\left[ \psi'(Y_i, \boldsymbol{\theta}_0) \right]$$
with $\psi'(Y_i, \boldsymbol{\theta}) = \frac{\partial}{\partial \boldsymbol{\theta}} \psi(Y_i, \boldsymbol{\theta})$. Otherwise,
$$\mathbf{A} = -\frac{\partial}{\partial \boldsymbol{\theta}} \mathbb{E}_g[\psi(Y_i, \boldsymbol{\theta})] \ \Biggr\lvert_{\boldsymbol{\theta} = \boldsymbol{\theta}_0}$$

\subsection*{PSD matrices}

Suppose that $\mathbf{x} \sim N(\boldsymbol{\mu}, \boldsymbol{\Sigma})$ is a multivariate normal random vector. We know that the covariance matrix $\boldsymbol{\Sigma}$ is positive semi-definite, and a useful fact about positive semi-definite matrices is that there exists a unique positive semi-definite matrix, which we will denote $\boldsymbol{\Sigma}^{\frac{1}{2}}$, such that $\boldsymbol{\Sigma}^{\frac{1}{2}}\boldsymbol{\Sigma}^{\frac{1}{2}} = \boldsymbol{\Sigma}$. It is also true that $\boldsymbol{\Sigma}^{-\frac{1}{2}} := (\boldsymbol{\Sigma}^{\frac{1}{2}})^{-1}$ satisfies $\boldsymbol{\Sigma}^{-\frac{1}{2}}\boldsymbol{\Sigma}^{-\frac{1}{2}} = \boldsymbol{\Sigma}^{-1}$.

\newpage

\section*{Questions}

\begin{enumerate}

\item Let $\{Y_n^{(1)}\}, \{Y_n^{(2)}\}, ..., \{Y_n^{(k)}\}$ be $k$ sequences of random variables, such that each $Y_n^{(i)} \overset{p}{\to} Y^{(i)}$ (here the superscript $(i)$ is used to distinguish the $i$th sequence; it does not denote an exponent, derivative, or order statistic). Show that
$$\max \limits_{i=1,...,k} |Y_n^{(i)} - Y^{(i)}| \overset{p}{\to} 0.$$

\bigskip

\item Let $X_1,...,X_n$ be iid from some distribution with cdf $F$. The \textit{empirical distribution function} $F_n$ is a \textit{nonparametric} estimate of $F$ defined by
$$F_n(t) = \frac{1}{n} \sum \limits_{i=1}^n \mathbbm{1}\{X_i \leq t \}.$$
On homework, you showed that for any $t$, $F_n(t) \overset{p}{\to} F(t)$. It turns out that we can make this result stronger. 

If $F$ is continuous, show that
$$\sup \limits_{t \in \mathbb{R}} |F_n(t) - F(t)| \overset{p}{\to} 0$$
(in other words, $F_n$ converges \textit{uniformly} to $F$ in probability; a slightly weaker version of the Glivenko-Cantelli theorem). \textit{Hint:} Begin by choosing $a_0, a_1,...,a_k$ such that 
$$-\infty = a_0 < a_1 < \cdots < a_k = \infty$$
and $F(a_i) - F(a_{i-1}) = \frac{1}{k}$ for $i = 1,...,k$ (your proof should explain why you can do this). You will also want to use the result of the previous problem, and you can use the homework result that $F_n(t) \overset{p}{\to} F(t)$ for each $t$ without proof.


\item Suppose that $(X_1, Y_1),...(X_n, Y_n)$ are iid observations from the \textbf{bivariate normal} distribution:

$$\begin{pmatrix}
X_i \\ Y_i
\end{pmatrix} \sim N\left( \begin{pmatrix}
\mu_1 \\ \mu_2
\end{pmatrix}, \ \begin{pmatrix}
\sigma_1^2 & \rho \sigma_1 \sigma_2 \\ 
\rho \sigma_1 \sigma_2 & \sigma_2^2
\end{pmatrix} \right)
$$

with $\mu_1, \mu_2 \in \mathbb{R}$, $\rho \in (-1, 1)$, and $\sigma_1, \sigma_2 > 0$. Find the method of moments estimators for $\mu_1, \mu_2, \sigma_1, \sigma_2$, and $\rho$.

\bigskip

\item Let $Y_1,...,Y_n$ be iid from some distribution with finite mean and variance $\mathbb{E}[Y_i] = \mu$ and $Var(Y_i) = \mu^4$. Find a variance stabilizing transformation $g$ such that the asymptotic variance of $\sqrt{n}(g(\overline{Y}) - g(\mu))$ does not depend on $\mu$.

\bigskip

\item Let $Y_1,...,Y_n \overset{iid}{\sim} Beta(\alpha, \beta)$, with pdf

$$f(y) = \frac{\Gamma(\alpha + \beta)}{\Gamma(\alpha) \Gamma(\beta)}y^{\alpha - 1} (1 - y)^{\beta - 1}$$

for $y \in (0, 1)$, and $\alpha, \beta > 0$. 

The mean is $\mathbb{E}[Y_i] = \dfrac{\alpha}{\alpha + \beta}$, and the variance is $Var(Y_i) = \dfrac{\alpha \beta}{(\alpha + \beta)^2(\alpha + \beta + 1)}$.

For this problem, it will help to define the \textit{digamma} function $\psi(u) = \frac{d}{du} \log \Gamma(u)$ and the \textit{trigamma} function $\psi^{(1)}(u) = \frac{d}{du} \psi(u)$. You may leave your answers in terms of these functions, unless otherwise specified. You may also use the following facts without proof:

$$\psi(u + 1) = \psi(u) + \frac{1}{u} \hspace{1cm} \psi^{(1)}(u + 1) = \psi^{(1)}(u) - \frac{1}{u^2}$$

\begin{enumerate}
\item Find the Fisher information matrix $\mathcal{I}_1(\alpha, \beta)$ for a single observation. You may assume that desired regularity conditions hold for the Beta distribution.

\item Now suppose that we \textit{know} the value of $\alpha$, and wish to estimate $\beta$. Find the method of moments estimator $\widehat{\beta}_{\text{MoM}}$.

\item Suppose that we furthermore \textit{know} that $\alpha = 1$. Find the MLE $\widehat{\beta}_{\text{MLE}}$. Your answer should be in closed form, and should not involve the digamma or trigamma functions when simplified.

\item Still supposing that we know $\alpha = 1$, find the asymptotic relative efficiency $ARE(\widehat{\beta}_{\text{MLE}}, \ \widehat{\beta}_{\text{MoM}})$. Your answer should be in closed form, and should not involve the digamma or trigamma functions when simplified.

\item How does the $ARE(\widehat{\beta}_{\text{MLE}}, \ \widehat{\beta}_{\text{MoM}})$ change as $\beta \to \infty$?

\item Explain the result from part (e) by discussing the behavior of the $Beta(1, \beta)$ distribution for large values of $\beta$, and the first-order Taylor expansion of the function $\log(1 - y)$ around 0.
\end{enumerate}




\bigskip

\item Let $Y_1,...,Y_n \overset{iid}{\sim} Exponential(\theta)$, with pdf 

$$f(y) = \theta e^{-\theta y} \hspace{1cm} y > 0$$

and cdf

$$F(y) = \begin{cases} 0 & y \leq 0 \\ 1 - e^{-\theta y} & y > 0 \end{cases}$$

The mean is $\mathbb{E}[Y] = 1/\theta$, and the variance is $Var(Y) = 1/\theta^2$.

\begin{enumerate}
\item Find the median of the $Exponential(\theta)$ distribution, as a function of $\theta$. (The population median, \textit{not} the sample median).

\item Let $M_n$ be the sample median computed from $Y_1,...,Y_n$. Using part (a), propose an estimator $\widehat{\theta}_{\text{med}}$ for $\theta$ that is a function of the sample median.

\item We have previously shown that the MLE for $\theta$ is $\widehat{\theta}_{\text{MLE}} = 1/\overline{Y}$. Calculate the asymptotic relative efficiency $ARE(\widehat{\theta}_{\text{MLE}}, \widehat{\theta}_{\text{med}})$.
\end{enumerate}

\item Suppose $X_1,...,X_n \overset{iid}{\sim} t_{\nu}$, with pdf

$$f(x) = \dfrac{\Gamma\left( \frac{\nu + 1}{2} \right)}{\sqrt{\pi \nu} \Gamma\left( \frac{\nu}{2} \right)} \left(1 + \frac{x^2}{\nu} \right)^{-\frac{(\nu + 1)}{2}}$$

The mean and median of the $t_\nu$ distribution are both 0, and the variance is $\dfrac{\nu}{\nu - 2}$ provided that $\nu > 2$.

Let 
$$Y_i = \mu + X_i$$
for some unknown $\mu \in \mathbb{R}$, and consider two estimators of $\mu$: the sample mean $\overline{Y}$, and the sample median $M_n$ calculated from $Y_1,...,Y_n$. Assuming $\nu > 2$, find the asymptotic relative efficiency $ARE(\overline{Y}, M_n)$. 

\bigskip

\item Let $Y_1,...,Y_n \overset{iid}{\sim} Poisson(\lambda)$, and let $\tau = P(Y_i = 0)$. We have previously shown that the MLE for $\lambda$ is $\widehat{\lambda}_{\text{MLE}} = \overline{Y}$.

\begin{enumerate}
\item Find the MLE $\widehat{\tau}_{\text{MLE}}$ for $\tau$.
\item Now consider another estimator: $\widetilde{\tau} = \frac{1}{n} \sum \limits_{i=1}^n \mathbbm{1}\{Y_i = 0 \}$. Find the asymptotic relative efficiency $ARE(\widehat{\tau}_{\text{MLE}}, \widetilde{\tau})$.
\end{enumerate}

\bigskip

\item Let $Y_1,...,Y_n \overset{iid}{\sim} N(\mu, \sigma^2)$. The median of the normal distribution is also $\mu$. The goal of this question is to compare two different estimators for the standard deviation $\sigma$:

$$\widehat{\sigma}_1 = \sqrt{\frac{1}{n} \sum \limits_{i=1}^n (Y_i - \overline{Y})^2} \hspace{2cm} \widehat{\sigma}_2 = \frac{\frac{1}{n} \sum \limits_{i=1}^n |Y_i - M_n|}{\sqrt{2/\pi}},$$

where $\overline{Y}$ is the sample mean and $M_n$ the sample median.

By itself, $\widehat{\sigma}_2$ is not an M-estimator, but it \textit{can} be made part of an M-estimator. Let $\boldsymbol{\theta} = (\theta_1, \theta_2)^T$ and let 
$$\widehat{\theta}_1 = M_n \hspace{1cm} \widehat{\theta}_2 = \widehat{\sigma}_2$$
Then, $\widehat{\boldsymbol{\theta}}$ solves
$$\sum \limits_{i=1}^n \psi(Y_i, \boldsymbol{\theta}) = \boldsymbol{c}_n \hspace{1cm} \text{where } \ \psi(Y_i, \boldsymbol{\theta}) = \begin{pmatrix}
\mathbbm{1}\{Y_i \leq \theta_1\} - \frac{1}{2} \\
\dfrac{|Y_i - \theta_1|}{\sqrt{2/\pi}} - \theta_2
\end{pmatrix}$$
and $\boldsymbol{c}_n/\sqrt{n} \overset{p}{\to} \boldsymbol{0}$. Let $\boldsymbol{\theta}_0$ solve $\mathbb{E}[\psi(Y_i, \boldsymbol{\theta})] = \boldsymbol{0}$.

\begin{enumerate}
\item Show that $\mathbb{E}\left(|Y_i - \mu|\right) = \sigma \sqrt{2/\pi}$, so $\widehat{\sigma}_2$ is a reasonable estimator for $\sigma$.

\item  Show that $\boldsymbol{\theta}_0 = (\mu, \sigma)^T$.

\item We have previously shown that $\sqrt{n}(\widehat{\sigma}_1^2 - \sigma^2) \overset{d}{\to} N(0, 2 \sigma^4)$. Using this information, find the asymptotic distribution of $\sqrt{n}(\widehat{\sigma}_1 - \sigma)$.

\item Using results for M-estimators, find the asymptotic distribution of $\sqrt{n}(\widehat{\boldsymbol{\theta}} - \boldsymbol{\theta}_0)$. For this part of the problem you may use, without proof, the following fact: 

$$\mathbb{E}(|Y_i - \theta_1|) = \sigma \sqrt{\frac{2}{\pi}} \ \exp \left\lbrace -\frac{1}{2\sigma^2} (\theta_1 - \mu)^2 \right\rbrace + (\mu - \theta_1)(1 - F(\theta_1)) + (\theta_1 - \mu)F(\theta_1)$$

where $F$ is the cdf of the $N(\mu, \sigma^2)$ distribution.

\item Find the asymptotic relative efficiency $ARE(\widehat{\sigma}_1, \widehat{\sigma}_2)$.

\end{enumerate}

\bigskip

\item Let $(\mathbf{x}_i, Y_i)$ be iid observations from some distribution $g$, and suppose that
$$Y_i = \exp\{\mathbf{x}_i^T \boldsymbol{\beta}\} + \varepsilon_i,$$
where $\mathbb{E}[\varepsilon_i] = 0$ and $Var(\varepsilon_i) = \sigma^2$. Our estimator $\widehat{\boldsymbol{\beta}}$ of the regression coefficients $\boldsymbol{\beta}$ minimizes

$$\sum \limits_{i=1}^n (Y_i - \exp\{\mathbf{x}_i^T \boldsymbol{\beta}\})^2$$

\begin{enumerate}
\item Find a function $\psi(Y_i, \mathbf{x}_i, \boldsymbol{\beta})$ such that $\widehat{\boldsymbol{\beta}}$ solves

$$\sum \limits_{i=1}^n \psi(Y_i, \mathbf{x}_i, \boldsymbol{\beta}) = \boldsymbol{0}$$

\item Let $\boldsymbol{\beta}^0$ be the value that solves $\mathbb{E}_g[\psi(Y_i, \mathbf{x}_i, \boldsymbol{\beta})] = \boldsymbol{0}$. Find the asymptotic distribution of $\sqrt{n}(\widehat{\boldsymbol{\beta}} - \boldsymbol{\beta}^0)$. The asymptotic variance should be in terms of $\boldsymbol{\beta}^0$, $\mathbf{x}_i$, and $Y_i$.

\item Give a plug-in estimator for $Var(\widehat{\boldsymbol{\beta}})$.
\end{enumerate}

\item The exponential distribution with parameter $\lambda$ has pdf
$$f(y | \lambda) = \frac{1}{\lambda} \exp\left\lbrace - \frac{y}{\lambda} \right\rbrace, \hspace{1cm} y > 0.$$
Let $Y_i > 0$ be a continuous, positive response variable of interest, and let $\mathbf{x}_i \in \mathbb{R}^{k+1}$ be a vector of covariates. Suppose we observe independent samples $(\mathbf{x}_1, Y_1),...,(\mathbf{x}_n, Y_n)$ from some distribution $g$, and we fit the following model (which may or may not be correct!):

\begin{align*}
Y_i|\mathbf{x}_i &\sim Exponential(\lambda_i) \\
-\frac{1}{\lambda_i} &= \mathbf{x}_i^T \boldsymbol{\beta}
\end{align*}
where the distribution of $\mathbf{x}_i$ does not depend on $\boldsymbol{\beta}$. 

\begin{enumerate}
\item \textbf{If our model is correct}, find the asymptotic distribution of $\sqrt{n}(\widehat{\boldsymbol{\beta}} - \boldsymbol{\beta})$, where $\boldsymbol{\beta}$ are the true regression coefficients in the population. Your answer for the asymptotic variance should be in terms of an expectation involving $\mathbf{x}_i$, $Y_i$, and $\boldsymbol{\beta}$. You may assume that desired regularity conditions hold for the Exponential.

\item Now suppose that our model is \textbf{mistaken}. Then, 
$$\sqrt{n}(\widehat{\boldsymbol{\beta}} - \boldsymbol{\beta}^0) \overset{d}{\to} N(\boldsymbol{0}, \mathbf{S}(\boldsymbol{\beta}^0))$$
Explain what $\boldsymbol{\beta}^0$ represents in this setting, in terms of some expectation with respect to the true distribution $g$. Then, find an expression for $\mathbf{S}(\boldsymbol{\beta}^0)$ in terms of expectations (with respect to $g$) involving $\mathbf{x}_i$, $Y_i$, and $\boldsymbol{\beta}^0$.

\item Give the plug-in estimator for $\mathbf{S}(\boldsymbol{\beta}^0)$ in terms of the observed $\mathbf{x}_i$, $Y_i$, and $\widehat{\boldsymbol{\beta}}$.
\end{enumerate}

\end{enumerate}

\end{document}
